% HD-R04. GPT-6 Astra Pro /root; 21 November 2026.
% Full proof; inherited effort not exposed.
\section{HD-R04: nonsimple fourfold secant seeds cannot give a semiregular product}
\label{sec:nonsimple-secant-general}

\paragraph{The exact scope.}
Let $X$ be an abelian fourfold, $H\in\operatorname{NS}(X)_{\mathbb Q}$
an ample class, $B\in\operatorname{NS}(X)_{\mathbb Q}$, and $s^2=-d$
with $d>0$. Fix a faithful finite point subgroup
$G\subset X\times\widehat X$ with an actual coherent translation/tensor
action, and an equivariant perfect complex $F$ with
$\operatorname{Ext}^{<0}(F,F)=0$ and
\[
 \operatorname{ch}(F)
   =e^B\bigl(\alpha e^{sH}+\bar\alpha e^{-sH}\bigr),
                   \qquad 0\ne\alpha\in\mathbb Q(s).       \tag{R04.1}
\]
Write $\mathcal C=D_{\mathrm{perf}}(X)^G$ and $\widetilde F$ for
its given equivariant object. At least one of the two categorical maps
\[
 \sigma^0_{\widetilde F}:\operatorname{End}_{\mathcal C}(\widetilde F)
                       \longrightarrow HH_0(\mathcal C),\qquad
 \sigma^2_{\widetilde F}:\operatorname{Ext}^2_{\mathcal C}
                   (\widetilde F,\widetilde F)
                       \longrightarrow HH_{-2}(\mathcal C)
                                                               \tag{R04.2}
\]
has a nonzero kernel. No simplicity or decomposition assumption is made.
In particular, the external product of any two objects (R04.1), with
its inherited product action, fails weak equivariant semiregularity.
Every actual Fourier--Mukai equivalence with the transported coherent
action retains that trace kernel. These assertions concern the stated
objects and actions; an unrelated action on the resulting eightfold
is not constructed here.

\paragraph{1. A real Clifford lower bound for every nonzero summand.}
For an abelian variety of dimension $g$, put
$V=H^1(X,\mathbb R)$ and $\mathcal W=(V\oplus V^*)_{\mathbb C}$,
with $Q(v+\xi,w+\eta)=\xi(w)+\eta(v)$ and
$c(v+\xi)=v\wedge+\iota_\xi$. The actual degree-one Hochschild
operators form
\[
 U=H^{0,1}(X)\oplus H^{1,0}(X)^*,\qquad
 \mathcal W=U\oplus\bar U,\qquad \dim_{\mathbb C}U=2g.
\]
Let $0\ne v\in H^{\mathrm{even}}(X,\mathbb R)$ be a real class,
and let $A=\{u\in U:c(u)v=0\}$. Complex conjugation gives
$c(\bar A)v=0$. Thus $A\oplus\bar A$ lies in the full Clifford
annihilator of $v$. This annihilator is isotropic: for any two of
its elements $a,b$, the anticommutator identity gives
\[
              Q(a,b)v=(c(a)c(b)+c(b)c(a))v=0.
\]
The form $Q$ is nondegenerate on the $4g$-dimensional space
$\mathcal W$. An isotropic subspace has dimension at most $2g$,
since it is contained in its perpendicular subspace. As
$U\cap\bar U=0$, this proves $2\dim A\le2g$. Consequently
\[
       \operatorname{rank}\bigl(U\longrightarrow H^*(X,\mathbb C),
                          \ u\longmapsto c(u)v\bigr)\ge g.       \tag{R04.3}
\]
For the rational Chern character of an actual perfect complex, this
map is precisely the trace of the degree-one Hochschild evaluation,
by the Atiyah identity used in HD-N11. The real form is the Betti
real form under de Rham comparison; the Chern character is real there.

The same lower bound applies to an equivariant direct summand whose
ordinary character is nonzero. Every point functor acts trivially on
$HH^1(X)$, and naturality makes its evaluation take values in the
invariant Ext group. Forgetting the equivariant structure carries its
trace to the ordinary trace. Thus, on a fourfold, every such summand
has invariant first self-Ext dimension at least four. This uses the
literal evaluation and forgetful maps, not a dimension comparison
between different moduli problems.

\paragraph{2. Degree-zero trace and the endomorphism algebra.}
For an actual perfect complex on a smooth variety, the degree-zero
Hochschild trace kills commutators and nilpotent endomorphisms.
Here are the exact maps needed for the latter assertion. Its HKR
components are
\[
 f\longmapsto\operatorname{tr}\!\left(
                         f\frac{\operatorname{At}(F)^p}{p!}\right)
                 \in H^p(X,\Omega_X^p).                         \tag{R04.4}
\]
These generalized traces are additive along a coherent cofiber
sequence of endomorphisms. The published locator is
Carmeli--Cnossen--Ramzi--Yanovski, \emph{Characters and transfer maps
via categorified traces}, Forum of Mathematics, Sigma\textbf{13}
(2025), e93, Theorem5.27, p.58, with the coherence requirement stated
in Remark5.25, p.57; DOI\href{https://doi.org/10.1017/fms.2025.23}
{10.1017/fms.2025.23}. Apply that generalized-trace statement in
$D_{\mathrm{qc}}(X)$ with source coefficient $\mathcal O_X$ and target
coefficient $\Omega_X^p[p]$. The Atiyah transformation is induced by
the universal diagonal kernel map; hence its powers are natural
transformations of exact derived functors. The functorial truncation
triangles, with their induced endomorphisms, supply the required
coherent cofiber diagrams for (R04.4).

For a nilpotent $f$ on $F$, these triangles express its trace as the
alternating sum of the traces of $H^j(f)$ on $H^j(F)$. Each $H^j(f)$
is an actual nilpotent sheaf map. Filter its sheaf by
$0\subset\ker H^j(f)\subset\ker H^j(f)^2\subset\cdots$.
The induced endomorphism on each successive quotient is zero.
All these sheaves are coherent and perfect because $X$ is smooth.
Their short exact sequences are coherent cofiber sequences, so the
same additivity proves that every component (R04.4) is zero.
This argument retains the sheaf filtration and makes no claim about
nilpotent traces in an arbitrary differential graded category.
Cyclicity and naturality of the Atiyah trace similarly give
$\sigma^0(fg)=\sigma^0(gf)$, proving the commutator assertion.

The point functors are exact for the coherent-sheaf $t$-structure,
so the preceding truncations and kernel filtrations are also
available equivariantly. Alternatively, forget them first. HD-N12's
actual isomorphism
\[
              HH_*(\mathcal C)\xrightarrow{\ \mathrm{Forg}_*\ }
                                      HH_*(X)                    \tag{R04.5}
\]
then reflects vanishing of each trace. Thus the same two assertions
hold for $\operatorname{End}_{\mathcal C}(\widetilde F)$.

Suppose now that the first map in (R04.2) is injective. Its
finite-dimensional endomorphism algebra is commutative, since every
commutator has zero trace. It has no nonzero nilpotent element by the
just-proved nilpotent calculation. A finite-dimensional commutative
reduced algebra over $\mathbb C$ is a product $\mathbb C^r$:
its finitely many maximal ideals have nilpotent intersection, which
is zero, and the Chinese remainder map identifies the product of
its residue fields. Idempotents split in the perfect equivariant
category, giving actual objects
\[
       \widetilde F=\bigoplus_{a=1}^rF_a,
       \qquad \operatorname{End}(F_a)=\mathbb C.                 \tag{R04.6}
\]
The trace of every identity idempotent is nonzero. By (R04.5) it
maps to the nonzero ordinary Chern character of its underlying
perfect complex. These characters are real rational Hodge classes.
Equation (R04.3) therefore applies to every $F_a$. Derived direct-sum
inclusions in degree one give
\[
        e_1:=\operatorname{ext}^1_{\mathcal C}
                       (\widetilde F,\widetilde F)\ge4r.        \tag{R04.7}
\]

\paragraph{3. The exact Euler contradiction.}
HD-N11 computes, for the entire character (R04.1),
\[
       \operatorname{rank}c^1=8,\qquad
       \operatorname{rank}c^2=12,\qquad
       \chi(F,F)=32d^2\alpha\bar\alpha\int_X H^4/4!>0.
                                                               \tag{R04.8}
\]
These are proved operator identities, including the $e^B$ shear.
HD-N12 identifies the weak equivariant trace with the ambient trace
on invariants, and proves
$\chi_{\mathcal C}(\widetilde F,\widetilde F)=\chi(F,F)/|G|>0$.
The degree-one evaluation gives $e_1\ge8$.
If the degree-two map in (R04.2) were injective as well, the invariant
Calabi--Yau fourfold Serre pairing would make its adjoint evaluation
surjective, by the literal pairing (N11.3). Hence
$\operatorname{ext}^2_{\mathcal C}(\widetilde F,\widetilde F)=12$.
Gluability and duality leave only Ext degrees zero through four.
Using (R04.6)--(R04.8), their alternating sum becomes
\[
 \begin{split}
 0<\chi_{\mathcal C}(\widetilde F,\widetilde F)
      &=2r-2e_1+12\\
      &\le\min\{2r-4,\,12-6r\}\le0.
 \end{split}                                             \tag{R04.9}
\]
The last inequality is exact for every positive integer $r$:
$r=1$ gives at most $-2$, $r=2$ gives at most zero, and $r\ge3$
gives at most $-6$. This contradiction proves the first assertion.

\paragraph{4. The actual external product and its transported kernel.}
Take two seeds as in (R04.1), on $X_1,X_2$, with the given coherent
groups $G_1,G_2$. In degree two their invariant external-product Ext
contains the direct Kunneth summands
\[
 \operatorname{Ext}^2_{G_1}(F_1,F_1)\otimes
       \operatorname{End}_{G_2}(F_2),\qquad
 \operatorname{End}_{G_1}(F_1)\otimes
       \operatorname{Ext}^2_{G_2}(F_2,F_2).                       \tag{R04.10}
\]
Their traces are the external products of the corresponding trace
maps; this follows from the external-product Atiyah identity used
in HD-N03. The degree-two trace--evaluation composite on either seed has rank
twelve. Evaluation therefore has rank at least twelve, and both
invariant second Ext spaces are nonzero.

A nonzero degree-two trace-kernel vector on one seed, tensored with
the identity of the other, is therefore a nonzero kernel vector in
(R04.10). A nonzero degree-zero trace-kernel vector, tensored with
any nonzero invariant second-Ext vector on the other seed, gives
another such nonzero vector. Paragraphs1--3 guarantee at least one
of these two alternatives on each seed. Thus the product always has
a nonzero invariant ambient trace kernel.

The forgetful Hochschild isomorphism (R04.5) applies on the product
as well: the nonidentity translation graphs are disjoint from the
diagonal, and nontrivial torsion tensor sectors have zero cohomology,
exactly as in HD-N12. This identifies the displayed kernel with a
kernel in the invariant category and proves failure of weak
$G_1\times G_2$-semiregularity. A specified Fourier--Mukai equivalence
carries (R04.10) isomorphically through its actual kernel functor,
and its Hochschild trace square commutes. Transporting the coherent
action therefore retains the nonzero kernel.

This closes the external-product secant-seed construction for arbitrary
nonsimple gluable fourfold seeds with the whole character (R04.1),
including the inherited faithful finite point actions. It extends
HD-N11--HD-N13 through the exact degree-zero and degree-two maps.
Other characters, longer coupled complexes and other actions remain
at their separately stated construction gates.
